% GENERATED FILE. Edit canonical lessons, then run npm run generate:books.
% canonical-source-hash: fnv1a64:44e1c83123accf55
% canonical-lessons: IT-C420-tabellone, IT-C420-ritardare, IT-C420-posticipare, IT-C420-rimborso, IT-C420-rimborsare, IT-R420-ripasso-413-420

\chapter{Ritardi e rimborsi}
\label{ch:it-ritardi-e-rimborsi}
\hlchaptermodality{420}

\begin{chapteropening}
\textbf{By the end of this chapter:} \emph{I can deal with delays: the departures board, running late, putting things back, a refund and refunding.}

Complete the last lesson of chapter 420: I can deal with delays: the departures board, running late, putting things back, a refund and refunding.
\end{chapteropening}

\section[il tabellone]{il tabellone — the departures board, the big board}

\label{lesson:IT-C420-tabellone}



\begin{quote}
Before the new one: say the Italian for the couchette, then the Italian for the guard.
\end{quote}

\subsection*{You'll want to know: il tabellone}
\textbf{il tabellone} — \textquotedblleft{}the departures board, the big board\textquotedblright{}.

\begin{quote}
Sul tabellone c'è scritto: ritardo venti minuti.
\end{quote}

\begin{quote}
\emph{The board says: twenty minutes late.}
\end{quote}

\emph{Il tabellone delle partenze} — departures; \emph{degli arrivi} — arrivals.

\begin{cousinweb}
\emph{tabellone} is a big \emph{tabella} 'board, chart', a learned borrowing from Latin \emph{tabella} 'small board', from \emph{tabula} 'board' (Lewis \& Short).
\end{cousinweb}

\begin{grammarlens}[title={What you've built}]
One more word for telling, explaining and describing at B1.
\end{grammarlens}

\subsection*{Your turn}
\begin{itemize}
\raggedright
  \item \emph{Say it:} \emph{il tabellone}
  \item \emph{Say it:} \emph{il tabellone}, once more
  \item \emph{Say it:} a sentence of your own with \emph{il tabellone}
  \item \emph{Recall:} say the Italian for the couchette, then the Italian for the guard, then say \emph{il tabellone} again
\end{itemize}

\begin{tcolorbox}[breakable,colback=teal!4,colframe=teal!35!black,title={Before you move on}]
What does il tabellone mean? (\textquotedblleft{}The departures board, the big board\textquotedblright{}.) Say it once more.
\end{tcolorbox}
\section[ritardare]{ritardare — to be late, to delay}

\label{lesson:IT-C420-ritardare}



\begin{quote}
Before the new one: say the Italian for the guard, then the Italian for the departures board.
\end{quote}

\subsection*{You'll want to know: ritardare}
\textbf{ritardare} — \textquotedblleft{}to be late, to delay\textquotedblright{}.

\begin{quote}
Il volo ritarda di due ore.
\end{quote}

\begin{quote}
\emph{The flight is two hours late.}
\end{quote}

\emph{In ritardo} — late.

\begin{cousinweb}
\emph{ritardare} is a learned borrowing from Latin \emph{retardāre} 'to hold back' (EtymDB), built on \emph{tardus} 'slow'.
\end{cousinweb}

\begin{grammarlens}[title={What you've built}]
One more word for telling, explaining and describing at B1.
\end{grammarlens}

\subsection*{Your turn}
\begin{itemize}
\raggedright
  \item \emph{Say it:} \emph{ritardare}
  \item \emph{Say it:} \emph{ritardare}, once more
  \item \emph{Say it:} a sentence of your own with \emph{ritardare}
  \item \emph{Recall:} say the Italian for the guard, then the Italian for the departures board, then say \emph{ritardare} again
\end{itemize}

\begin{tcolorbox}[breakable,colback=teal!4,colframe=teal!35!black,title={Before you move on}]
What does ritardare mean? (\textquotedblleft{}To be late, to delay\textquotedblright{}.) Say it once more.
\end{tcolorbox}
\section[posticipare]{posticipare — to postpone, to put back}

\label{lesson:IT-C420-posticipare}



\begin{quote}
Before the new one: say the Italian for the departures board, then the Italian for to be late.
\end{quote}

\subsection*{You'll want to know: posticipare}
\textbf{posticipare} — \textquotedblleft{}to postpone, to put back\textquotedblright{}.

\begin{quote}
Hanno posticipato la partenza a domani.
\end{quote}

\begin{quote}
\emph{They put the departure back to tomorrow.}
\end{quote}

A more formal word than \emph{rimandare}.

\begin{cousinweb}
\emph{posticipare} is a learned formation on Latin \emph{post} 'after' (Lewis \& Short), shaped on the model of \emph{anticipare} 'to bring forward'.
\end{cousinweb}

\begin{grammarlens}[title={What you've built}]
One more word for telling, explaining and describing at B1.
\end{grammarlens}

\subsection*{Your turn}
\begin{itemize}
\raggedright
  \item \emph{Say it:} \emph{posticipare}
  \item \emph{Say it:} \emph{posticipare}, once more
  \item \emph{Say it:} a sentence of your own with \emph{posticipare}
  \item \emph{Recall:} say the Italian for the departures board, then the Italian for to be late, then say \emph{posticipare} again
\end{itemize}

\begin{tcolorbox}[breakable,colback=teal!4,colframe=teal!35!black,title={Before you move on}]
What does posticipare mean? (\textquotedblleft{}To postpone, to put back\textquotedblright{}.) Say it once more.
\end{tcolorbox}
\section[il rimborso]{il rimborso — the refund}

\label{lesson:IT-C420-rimborso}



\begin{quote}
Before the new one: say the Italian for to be late, then the Italian for to postpone.
\end{quote}

\subsection*{You'll want to know: il rimborso}
\textbf{il rimborso} — \textquotedblleft{}the refund\textquotedblright{}.

\begin{quote}
Ho chiesto il rimborso del biglietto.
\end{quote}

\begin{quote}
\emph{I asked for a refund on the ticket.}
\end{quote}

\emph{Un rimborso parziale} — a partial refund.

\begin{cousinweb}
\emph{rimborso} is built on \emph{borsa} 'purse', which is inherited from Late Latin \emph{bursa}, from Greek \emph{býrsa} 'hide, leather': money back into the purse.
\end{cousinweb}

\begin{grammarlens}[title={What you've built}]
One more word for telling, explaining and describing at B1.
\end{grammarlens}

\subsection*{Your turn}
\begin{itemize}
\raggedright
  \item \emph{Say it:} \emph{il rimborso}
  \item \emph{Say it:} \emph{il rimborso}, once more
  \item \emph{Say it:} a sentence of your own with \emph{il rimborso}
  \item \emph{Recall:} say the Italian for to be late, then the Italian for to postpone, then say \emph{il rimborso} again
\end{itemize}

\begin{tcolorbox}[breakable,colback=teal!4,colframe=teal!35!black,title={Before you move on}]
What does il rimborso mean? (\textquotedblleft{}The refund\textquotedblright{}.) Say it once more.
\end{tcolorbox}
\section[rimborsare]{rimborsare — to refund, to pay back}

\label{lesson:IT-C420-rimborsare}



\begin{quote}
Before the new one: say the Italian for to postpone, then the Italian for the refund.
\end{quote}

\subsection*{You'll want to know: rimborsare}
\textbf{rimborsare} — \textquotedblleft{}to refund, to pay back\textquotedblright{}.

\begin{quote}
Ci hanno rimborsato metà del prezzo.
\end{quote}

\begin{quote}
\emph{They refunded us half the price.}
\end{quote}

\emph{Farsi rimborsare} — to get a refund.

\begin{cousinweb}
\emph{rimborsare} is built from \emph{ri-} 'back' and \emph{borsa} 'purse', which is inherited from Late Latin \emph{bursa}, from Greek \emph{býrsa} 'hide, leather'.
\end{cousinweb}

\begin{grammarlens}[title={What you've built}]
One more word for telling, explaining and describing at B1.
\end{grammarlens}

\subsection*{Your turn}
\begin{itemize}
\raggedright
  \item \emph{Say it:} \emph{rimborsare}
  \item \emph{Say it:} \emph{rimborsare}, once more
  \item \emph{Say it:} a sentence of your own with \emph{rimborsare}
  \item \emph{Recall:} say the Italian for to postpone, then the Italian for the refund, then say \emph{rimborsare} again
\end{itemize}

\begin{tcolorbox}[breakable,colback=teal!4,colframe=teal!35!black,title={Before you move on}]
What does rimborsare mean? (\textquotedblleft{}To refund, to pay back\textquotedblright{}.) Say it once more.
\end{tcolorbox}
\section[(dialogue)]{Ripasso dei capitoli 413–420}

\label{lesson:IT-R420-ripasso-413-420}



\begin{quote}
Say the Italian for to rear-end and the rear-end collision.
\end{quote}

\subsection*{Your turn}
Say these aloud:

\begin{itemize}
\raggedright
  \item to rear-end, the rear-end collision, the traffic jam, the diversion, to divert
  \item the main road, the provincial road, the bend, the pothole, the dirt road
  \item to hire, the hire, the number plate, the logbook, the insurance
  \item the moped, the scooter, to pedal, the parking meter, to stop briefly
  \item the platform, the bus shelter, the terminus, the long-distance bus, to pass through
  \item the ticket office, the season ticket, to stamp, the ticket inspector, the supplement
  \item the railway carriage, the coach, the compartment, the couchette, the guard
  \item the departures board, to be late, to postpone, the refund, to refund
\end{itemize}

\begin{tcolorbox}[breakable,colback=teal!4,colframe=teal!35!black,title={Before you move on}]
To rear-end? (\textbf{tamponare}.) The rear-end collision? (\textbf{il tamponamento}.) The traffic jam? (\textbf{l'ingorgo}.) The diversion? (\textbf{la deviazione}.) To divert? (\textbf{deviare}.) The main road? (\textbf{la statale}.) The provincial road? (\textbf{la provinciale}.) The bend? (\textbf{la curva}.) The pothole? (\textbf{la buca}.) The dirt road? (\textbf{lo sterrato}.) To hire? (\textbf{noleggiare}.) The hire? (\textbf{il noleggio}.) The number plate? (\textbf{la targa}.) The logbook? (\textbf{il libretto}.) The insurance? (\textbf{l'assicurazione}.) The moped? (\textbf{il motorino}.) The scooter? (\textbf{lo scooter}.) To pedal? (\textbf{pedalare}.) The parking meter? (\textbf{il parchimetro}.) To stop briefly? (\textbf{sostare}.) The platform? (\textbf{la banchina}.) The bus shelter? (\textbf{la pensilina}.) The terminus? (\textbf{il capolinea}.) The long-distance bus? (\textbf{il pullman}.) To pass through? (\textbf{transitare}.) The ticket office? (\textbf{la biglietteria}.) The season ticket? (\textbf{l'abbonamento}.) To stamp? (\textbf{timbrare}.) The ticket inspector? (\textbf{il controllore}.) The supplement? (\textbf{il supplemento}.) The railway carriage? (\textbf{il vagone}.) The coach? (\textbf{la carrozza}.) The compartment? (\textbf{lo scompartimento}.) The couchette? (\textbf{la cuccetta}.) The guard? (\textbf{il capotreno}.) The departures board? (\textbf{il tabellone}.) To be late? (\textbf{ritardare}.) To postpone? (\textbf{posticipare}.) The refund? (\textbf{il rimborso}.) To refund? (\textbf{rimborsare}.)
\end{tcolorbox}
